\documentclass[12pt, a4paper]{report}

% ── packages ──────────────────────────────────────────────────────────────────
\usepackage[T1]{fontenc}
\usepackage[utf8]{inputenc}
\usepackage[a4paper, margin=2.5cm]{geometry}
\usepackage{graphicx}
\usepackage{booktabs}
\usepackage{amsmath}
\usepackage{amssymb}
\usepackage{xcolor}
\usepackage{hyperref}
\usepackage{microtype}
\usepackage{siunitx}
\usepackage{subcaption}
\usepackage{setspace}
\usepackage{fancyhdr}
\usepackage[acronym]{glossaries}

\sisetup{detect-weight=true, mode=text}
\onehalfspacing

% ── figure paths ──────────────────────────────────────────────────────────────
\graphicspath{{../../results/}}

% ── convenience ───────────────────────────────────────────────────────────────
\newcommand{\tbd}{\textcolor{gray}{--}}
\newcommand{\na}{\textcolor{gray}{n/a}}

% ── graceful fallbacks for auto-generated metrics ─────────────────────────────
\providecommand{\whuGAITCnnTestAcc}{\tbd}
\providecommand{\whuGAITAuthTestAUC}{\tbd}
\providecommand{\whuGAITSimplDAUC}{\tbd}
\providecommand{\whuGAITComboASR}{\tbd}
\providecommand{\whuGAITBudgetRatio}{\tbd}
\providecommand{\uciharCnnTestAcc}{\tbd}
\providecommand{\uciharAuthTestAUC}{\tbd}
\providecommand{\uciharSimplDAUC}{\tbd}
\providecommand{\uciharComboASR}{\tbd}
\providecommand{\wisdmCnnTestAcc}{\tbd}
\providecommand{\wisdmAuthTestAUC}{\tbd}
\providecommand{\wisdmSimplDAUC}{\tbd}
\providecommand{\wisdmComboASR}{\tbd}
\providecommand{\combinedAuthTestAUC}{\tbd}
\providecommand{\combinedSimplDAUC}{\tbd}
\providecommand{\signalAuthModelAUCB}{\tbd}
\providecommand{\signalAuthModelAUCA}{\tbd}
\providecommand{\signalAuthModelAUCC}{\tbd}
\providecommand{\signalMiaAUCBvsC}{\tbd}
\providecommand{\signalMiaAUCAvsC}{\tbd}

% ── auto-generated metrics ────────────────────────────────────────────────────
\IfFileExists{../generated/compare_models_whuGAIT_metrics.tex}
  {% Auto-generated by compare_models_whuGAIT: DO NOT EDIT MANUALLY
% Last run: 2026-07-28 13:26:26

\providecommand{\whuGAITCnnTestAcc}{}
\renewcommand{\whuGAITCnnTestAcc}{91.44\%}
\providecommand{\whuGAITAuthTestAcc}{}
\renewcommand{\whuGAITAuthTestAcc}{80.12\%}
\providecommand{\whuGAITAuthTrainAcc}{}
\renewcommand{\whuGAITAuthTrainAcc}{91.54\%}
\providecommand{\whuGAITAuthTestAUC}{}
\renewcommand{\whuGAITAuthTestAUC}{0.8698}
\providecommand{\whuGAITAuthOverfitAcc}{}
\renewcommand{\whuGAITAuthOverfitAcc}{+11.42pp}
\providecommand{\whuGAITAuthSameMean}{}
\renewcommand{\whuGAITAuthSameMean}{0.1480}
\providecommand{\whuGAITAuthDiffMean}{}
\renewcommand{\whuGAITAuthDiffMean}{0.6710}
\providecommand{\whuGAITAuthSep}{}
\renewcommand{\whuGAITAuthSep}{0.5220}
}{}
\IfFileExists{../generated/compare_models_ucihar_metrics.tex}
  {% Auto-generated by compare_models_ucihar: DO NOT EDIT MANUALLY
% Last run: 2026-07-28 13:26:26

\providecommand{\uciharCnnTestAcc}{}
\renewcommand{\uciharCnnTestAcc}{85.34\%}
\providecommand{\uciharAuthTestAcc}{}
\renewcommand{\uciharAuthTestAcc}{68.83\%}
\providecommand{\uciharAuthTrainAcc}{}
\renewcommand{\uciharAuthTrainAcc}{81.61\%}
\providecommand{\uciharAuthTestAUC}{}
\renewcommand{\uciharAuthTestAUC}{0.7388}
\providecommand{\uciharAuthOverfitAcc}{}
\renewcommand{\uciharAuthOverfitAcc}{+12.78pp}
\providecommand{\uciharAuthSameMean}{}
\renewcommand{\uciharAuthSameMean}{0.4290}
\providecommand{\uciharAuthDiffMean}{}
\renewcommand{\uciharAuthDiffMean}{0.7200}
\providecommand{\uciharAuthSep}{}
\renewcommand{\uciharAuthSep}{0.2900}
}{}
\IfFileExists{../generated/compare_models_wisdm_metrics.tex}
  {% Auto-generated by compare_models_wisdm: DO NOT EDIT MANUALLY
% Last run: 2026-07-28 13:26:26

\providecommand{\wisdmCnnTestAcc}{}
\renewcommand{\wisdmCnnTestAcc}{96.29\%}
\providecommand{\wisdmAuthTestAcc}{}
\renewcommand{\wisdmAuthTestAcc}{91.87\%}
\providecommand{\wisdmAuthTrainAcc}{}
\renewcommand{\wisdmAuthTrainAcc}{98.02\%}
\providecommand{\wisdmAuthTestAUC}{}
\renewcommand{\wisdmAuthTestAUC}{0.9577}
\providecommand{\wisdmAuthOverfitAcc}{}
\renewcommand{\wisdmAuthOverfitAcc}{+6.15pp}
\providecommand{\wisdmAuthSameMean}{}
\renewcommand{\wisdmAuthSameMean}{0.0870}
\providecommand{\wisdmAuthDiffMean}{}
\renewcommand{\wisdmAuthDiffMean}{0.8790}
\providecommand{\wisdmAuthSep}{}
\renewcommand{\wisdmAuthSep}{0.7920}
}{}
\IfFileExists{../generated/compare_models_combined_metrics.tex}
  {% Auto-generated by compare_models_combined: DO NOT EDIT MANUALLY
% Last run: 2026-08-06 21:44:47

\providecommand{\combinedCnnTestAcc}{}
\renewcommand{\combinedCnnTestAcc}{—}
\providecommand{\combinedAuthTestAcc}{}
\renewcommand{\combinedAuthTestAcc}{75.72\%}
\providecommand{\combinedAuthTrainAcc}{}
\renewcommand{\combinedAuthTrainAcc}{82.95\%}
\providecommand{\combinedAuthTestAUC}{}
\renewcommand{\combinedAuthTestAUC}{0.8378}
\providecommand{\combinedAuthOverfitAcc}{}
\renewcommand{\combinedAuthOverfitAcc}{+7.23pp}
\providecommand{\combinedAuthSameMean}{}
\renewcommand{\combinedAuthSameMean}{0.3430}
\providecommand{\combinedAuthDiffMean}{}
\renewcommand{\combinedAuthDiffMean}{0.7370}
\providecommand{\combinedAuthSep}{}
\renewcommand{\combinedAuthSep}{0.3940}
}{}
\IfFileExists{../generated/compare_mia_whuGAIT_metrics.tex}
  {% Auto-generated by compare_mia_whuGAIT: DO NOT EDIT MANUALLY
% Last run: 2026-09-08 12:25:56

\providecommand{\whuGAITSimplDAUC}{}
\renewcommand{\whuGAITSimplDAUC}{0.815}
\providecommand{\whuGAITSimplDDeltaGap}{}
\renewcommand{\whuGAITSimplDDeltaGap}{0.233}
\providecommand{\whuGAITSimplDTPRten}{}
\renewcommand{\whuGAITSimplDTPRten}{0.298}
\providecommand{\whuGAITGreyLiraRawAUC}{}
\renewcommand{\whuGAITGreyLiraRawAUC}{0.262}
\providecommand{\whuGAITGreyLiraRawACC}{}
\renewcommand{\whuGAITGreyLiraRawACC}{0.183}
\providecommand{\whuGAITGreyLiraRawTPRten}{}
\renewcommand{\whuGAITGreyLiraRawTPRten}{0.000}
\providecommand{\whuGAITGreyLiraLfAUC}{}
\renewcommand{\whuGAITGreyLiraLfAUC}{0.163}
\providecommand{\whuGAITGreyLiraLfACC}{}
\renewcommand{\whuGAITGreyLiraLfACC}{0.183}
\providecommand{\whuGAITGreyLiraLfTPRten}{}
\renewcommand{\whuGAITGreyLiraLfTPRten}{0.000}
\providecommand{\whuGAITGreyLiraMfAUC}{}
\renewcommand{\whuGAITGreyLiraMfAUC}{0.251}
\providecommand{\whuGAITGreyLiraMfACC}{}
\renewcommand{\whuGAITGreyLiraMfACC}{0.183}
\providecommand{\whuGAITGreyLiraMfTPRten}{}
\renewcommand{\whuGAITGreyLiraMfTPRten}{0.000}
\providecommand{\whuGAITWarmLiraRawAUC}{}
\renewcommand{\whuGAITWarmLiraRawAUC}{0.346}
\providecommand{\whuGAITWarmLiraRawACC}{}
\renewcommand{\whuGAITWarmLiraRawACC}{0.394}
\providecommand{\whuGAITWarmLiraRawTPRten}{}
\renewcommand{\whuGAITWarmLiraRawTPRten}{0.048}
\providecommand{\whuGAITColdLiraRawAUC}{}
\renewcommand{\whuGAITColdLiraRawAUC}{0.709}
\providecommand{\whuGAITColdLiraRawACC}{}
\renewcommand{\whuGAITColdLiraRawACC}{0.817}
\providecommand{\whuGAITColdLiraRawTPRten}{}
\renewcommand{\whuGAITColdLiraRawTPRten}{0.417}
\providecommand{\whuGAITColdLiraLfAUC}{}
\renewcommand{\whuGAITColdLiraLfAUC}{0.686}
\providecommand{\whuGAITColdLiraLfACC}{}
\renewcommand{\whuGAITColdLiraLfACC}{0.808}
\providecommand{\whuGAITColdLiraLfTPRten}{}
\renewcommand{\whuGAITColdLiraLfTPRten}{0.381}
\providecommand{\whuGAITColdLiraMfAUC}{}
\renewcommand{\whuGAITColdLiraMfAUC}{0.743}
\providecommand{\whuGAITColdLiraMfACC}{}
\renewcommand{\whuGAITColdLiraMfACC}{0.827}
\providecommand{\whuGAITColdLiraMfTPRten}{}
\renewcommand{\whuGAITColdLiraMfTPRten}{0.274}
}{}
\IfFileExists{../generated/compare_mia_ucihar_metrics.tex}
  {% Auto-generated by compare_mia_ucihar: DO NOT EDIT MANUALLY
% Last run: 2026-09-08 12:25:56

\providecommand{\uciharSimplDAUC}{}
\renewcommand{\uciharSimplDAUC}{0.794}
\providecommand{\uciharSimplDDeltaGap}{}
\renewcommand{\uciharSimplDDeltaGap}{0.188}
\providecommand{\uciharSimplDTPRten}{}
\renewcommand{\uciharSimplDTPRten}{0.333}
\providecommand{\uciharGreyLiraRawAUC}{}
\renewcommand{\uciharGreyLiraRawAUC}{0.889}
\providecommand{\uciharGreyLiraRawACC}{}
\renewcommand{\uciharGreyLiraRawACC}{0.800}
\providecommand{\uciharGreyLiraRawTPRten}{}
\renewcommand{\uciharGreyLiraRawTPRten}{0.714}
\providecommand{\uciharGreyLiraLfAUC}{}
\renewcommand{\uciharGreyLiraLfAUC}{0.815}
\providecommand{\uciharGreyLiraLfACC}{}
\renewcommand{\uciharGreyLiraLfACC}{0.767}
\providecommand{\uciharGreyLiraLfTPRten}{}
\renewcommand{\uciharGreyLiraLfTPRten}{0.667}
\providecommand{\uciharGreyLiraMfAUC}{}
\renewcommand{\uciharGreyLiraMfAUC}{0.862}
\providecommand{\uciharGreyLiraMfACC}{}
\renewcommand{\uciharGreyLiraMfACC}{0.700}
\providecommand{\uciharGreyLiraMfTPRten}{}
\renewcommand{\uciharGreyLiraMfTPRten}{0.571}
\providecommand{\uciharWarmLiraRawAUC}{}
\renewcommand{\uciharWarmLiraRawAUC}{0.661}
\providecommand{\uciharWarmLiraRawACC}{}
\renewcommand{\uciharWarmLiraRawACC}{0.667}
\providecommand{\uciharWarmLiraRawTPRten}{}
\renewcommand{\uciharWarmLiraRawTPRten}{0.000}
\providecommand{\uciharColdLiraRawAUC}{}
\renewcommand{\uciharColdLiraRawAUC}{0.968}
\providecommand{\uciharColdLiraRawACC}{}
\renewcommand{\uciharColdLiraRawACC}{0.700}
\providecommand{\uciharColdLiraRawTPRten}{}
\renewcommand{\uciharColdLiraRawTPRten}{0.905}
\providecommand{\uciharColdLiraLfAUC}{}
\renewcommand{\uciharColdLiraLfAUC}{0.921}
\providecommand{\uciharColdLiraLfACC}{}
\renewcommand{\uciharColdLiraLfACC}{0.700}
\providecommand{\uciharColdLiraLfTPRten}{}
\renewcommand{\uciharColdLiraLfTPRten}{0.667}
\providecommand{\uciharColdLiraMfAUC}{}
\renewcommand{\uciharColdLiraMfAUC}{0.926}
\providecommand{\uciharColdLiraMfACC}{}
\renewcommand{\uciharColdLiraMfACC}{0.700}
\providecommand{\uciharColdLiraMfTPRten}{}
\renewcommand{\uciharColdLiraMfTPRten}{0.762}
}{}
\IfFileExists{../generated/compare_mia_wisdm_metrics.tex}
  {% Auto-generated by compare_mia_wisdm: DO NOT EDIT MANUALLY
% Last run: 2026-09-08 12:25:56

\providecommand{\wisdmSimplDAUC}{}
\renewcommand{\wisdmSimplDAUC}{0.952}
\providecommand{\wisdmSimplDDeltaGap}{}
\renewcommand{\wisdmSimplDDeltaGap}{0.136}
\providecommand{\wisdmSimplDTPRten}{}
\renewcommand{\wisdmSimplDTPRten}{0.950}
\providecommand{\wisdmGreyLiraRawAUC}{}
\renewcommand{\wisdmGreyLiraRawAUC}{0.827}
\providecommand{\wisdmGreyLiraRawACC}{}
\renewcommand{\wisdmGreyLiraRawACC}{0.706}
\providecommand{\wisdmGreyLiraRawTPRten}{}
\renewcommand{\wisdmGreyLiraRawTPRten}{0.625}
\providecommand{\wisdmGreyLiraLfAUC}{}
\renewcommand{\wisdmGreyLiraLfAUC}{0.723}
\providecommand{\wisdmGreyLiraLfACC}{}
\renewcommand{\wisdmGreyLiraLfACC}{0.608}
\providecommand{\wisdmGreyLiraLfTPRten}{}
\renewcommand{\wisdmGreyLiraLfTPRten}{0.525}
\providecommand{\wisdmGreyLiraMfAUC}{}
\renewcommand{\wisdmGreyLiraMfAUC}{0.857}
\providecommand{\wisdmGreyLiraMfACC}{}
\renewcommand{\wisdmGreyLiraMfACC}{0.765}
\providecommand{\wisdmGreyLiraMfTPRten}{}
\renewcommand{\wisdmGreyLiraMfTPRten}{0.725}
\providecommand{\wisdmWarmLiraRawAUC}{}
\renewcommand{\wisdmWarmLiraRawAUC}{0.659}
\providecommand{\wisdmWarmLiraRawACC}{}
\renewcommand{\wisdmWarmLiraRawACC}{0.784}
\providecommand{\wisdmWarmLiraRawTPRten}{}
\renewcommand{\wisdmWarmLiraRawTPRten}{0.250}
\providecommand{\wisdmColdLiraRawAUC}{}
\renewcommand{\wisdmColdLiraRawAUC}{0.916}
\providecommand{\wisdmColdLiraRawACC}{}
\renewcommand{\wisdmColdLiraRawACC}{0.784}
\providecommand{\wisdmColdLiraRawTPRten}{}
\renewcommand{\wisdmColdLiraRawTPRten}{0.800}
\providecommand{\wisdmColdLiraLfAUC}{}
\renewcommand{\wisdmColdLiraLfAUC}{0.968}
\providecommand{\wisdmColdLiraLfACC}{}
\renewcommand{\wisdmColdLiraLfACC}{0.784}
\providecommand{\wisdmColdLiraLfTPRten}{}
\renewcommand{\wisdmColdLiraLfTPRten}{0.925}
\providecommand{\wisdmColdLiraMfAUC}{}
\renewcommand{\wisdmColdLiraMfAUC}{0.968}
\providecommand{\wisdmColdLiraMfACC}{}
\renewcommand{\wisdmColdLiraMfACC}{0.804}
\providecommand{\wisdmColdLiraMfTPRten}{}
\renewcommand{\wisdmColdLiraMfTPRten}{0.950}
}{}
\IfFileExists{../generated/compare_mia_combined_metrics.tex}
  {% Auto-generated by compare_mia_combined: DO NOT EDIT MANUALLY
% Last run: 2026-07-28 13:55:15

\providecommand{\combinedSimplDAUC}{}
\renewcommand{\combinedSimplDAUC}{0.753}
\providecommand{\combinedSimplDDeltaGap}{}
\renewcommand{\combinedSimplDDeltaGap}{0.262}
\providecommand{\combinedGreyLiraRawACC}{}
\renewcommand{\combinedGreyLiraRawACC}{—}
\providecommand{\combinedGreyLiraRawAUC}{}
\renewcommand{\combinedGreyLiraRawAUC}{—}
\providecommand{\combinedGreyLiraLfACC}{}
\renewcommand{\combinedGreyLiraLfACC}{—}
\providecommand{\combinedGreyLiraLfAUC}{}
\renewcommand{\combinedGreyLiraLfAUC}{—}
\providecommand{\combinedGreyLiraMfACC}{}
\renewcommand{\combinedGreyLiraMfACC}{—}
\providecommand{\combinedGreyLiraMfAUC}{}
\renewcommand{\combinedGreyLiraMfAUC}{—}
\providecommand{\combinedWarmLiraRawACC}{}
\renewcommand{\combinedWarmLiraRawACC}{—}
\providecommand{\combinedWarmLiraRawAUC}{}
\renewcommand{\combinedWarmLiraRawAUC}{—}
\providecommand{\combinedWarmLiraLfACC}{}
\renewcommand{\combinedWarmLiraLfACC}{—}
\providecommand{\combinedWarmLiraLfAUC}{}
\renewcommand{\combinedWarmLiraLfAUC}{—}
\providecommand{\combinedWarmLiraMfACC}{}
\renewcommand{\combinedWarmLiraMfACC}{—}
\providecommand{\combinedWarmLiraMfAUC}{}
\renewcommand{\combinedWarmLiraMfAUC}{—}
\providecommand{\combinedColdLiraRawACC}{}
\renewcommand{\combinedColdLiraRawACC}{—}
\providecommand{\combinedColdLiraRawAUC}{}
\renewcommand{\combinedColdLiraRawAUC}{—}
\providecommand{\combinedColdLiraLfACC}{}
\renewcommand{\combinedColdLiraLfACC}{—}
\providecommand{\combinedColdLiraLfAUC}{}
\renewcommand{\combinedColdLiraLfAUC}{—}
\providecommand{\combinedColdLiraMfACC}{}
\renewcommand{\combinedColdLiraMfACC}{—}
\providecommand{\combinedColdLiraMfAUC}{}
\renewcommand{\combinedColdLiraMfAUC}{—}
}{}
\IfFileExists{../generated/compare_evasion_whuGAIT_metrics.tex}
  {% Auto-generated by compare_evasion_whuGAIT: DO NOT EDIT MANUALLY
% Last run: 2026-08-06 21:45:02

\providecommand{\whuGAITEpsTarget}{}
\renewcommand{\whuGAITEpsTarget}{24.54}
\providecommand{\whuGAITComboASR}{}
\renewcommand{\whuGAITComboASR}{1.000}
\providecommand{\whuGAITBudgetRatio}{}
\renewcommand{\whuGAITBudgetRatio}{0.056}
\providecommand{\whuGAITNResistant}{}
\renewcommand{\whuGAITNResistant}{0}
\providecommand{\whuGAITNPairs}{}
\renewcommand{\whuGAITNPairs}{3800}
\providecommand{\whuGAITNFullCombos}{}
\renewcommand{\whuGAITNFullCombos}{220}
\providecommand{\whuGAITNCombos}{}
\renewcommand{\whuGAITNCombos}{220}
\providecommand{\whuGAITTrainComboASR}{}
\renewcommand{\whuGAITTrainComboASR}{0.914}
\providecommand{\whuGAITTrainBudgetRatio}{}
\renewcommand{\whuGAITTrainBudgetRatio}{0.269}
}{}
\IfFileExists{../generated/compare_evasion_ucihar_metrics.tex}
  {% Auto-generated by compare_evasion_ucihar: DO NOT EDIT MANUALLY
% Last run: 2026-07-28 13:26:39

\providecommand{\uciharEpsTarget}{}
\renewcommand{\uciharEpsTarget}{37.47}
\providecommand{\uciharComboASR}{}
\renewcommand{\uciharComboASR}{0.980}
\providecommand{\uciharBudgetRatio}{}
\renewcommand{\uciharBudgetRatio}{0.135}
\providecommand{\uciharNResistant}{}
\renewcommand{\uciharNResistant}{106}
\providecommand{\uciharNPairs}{}
\renewcommand{\uciharNPairs}{5000}
\providecommand{\uciharNFullCombos}{}
\renewcommand{\uciharNFullCombos}{42}
\providecommand{\uciharNCombos}{}
\renewcommand{\uciharNCombos}{72}
\providecommand{\uciharTrainComboASR}{}
\renewcommand{\uciharTrainComboASR}{0.970}
\providecommand{\uciharTrainBudgetRatio}{}
\renewcommand{\uciharTrainBudgetRatio}{0.148}
}{}
\IfFileExists{../generated/compare_evasion_wisdm_metrics.tex}
  {% Auto-generated by compare_evasion_wisdm: DO NOT EDIT MANUALLY
% Last run: 2026-07-28 13:26:39

\providecommand{\wisdmEpsTarget}{}
\renewcommand{\wisdmEpsTarget}{31.00}
\providecommand{\wisdmComboASR}{}
\renewcommand{\wisdmComboASR}{0.978}
\providecommand{\wisdmBudgetRatio}{}
\renewcommand{\wisdmBudgetRatio}{0.300}
\providecommand{\wisdmNResistant}{}
\renewcommand{\wisdmNResistant}{106}
\providecommand{\wisdmNPairs}{}
\renewcommand{\wisdmNPairs}{5000}
\providecommand{\wisdmNFullCombos}{}
\renewcommand{\wisdmNFullCombos}{88}
\providecommand{\wisdmNCombos}{}
\renewcommand{\wisdmNCombos}{110}
\providecommand{\wisdmTrainComboASR}{}
\renewcommand{\wisdmTrainComboASR}{0.932}
\providecommand{\wisdmTrainBudgetRatio}{}
\renewcommand{\wisdmTrainBudgetRatio}{0.415}
}{}
\IfFileExists{../generated/compare_evasion_combined_metrics.tex}
  {% Auto-generated by compare_evasion_combined: DO NOT EDIT MANUALLY
% Last run: 2026-07-28 13:26:39

\providecommand{\combinedEpsTarget}{}
\renewcommand{\combinedEpsTarget}{24.11}
\providecommand{\combinedComboASR}{}
\renewcommand{\combinedComboASR}{0.962}
\providecommand{\combinedBudgetRatio}{}
\renewcommand{\combinedBudgetRatio}{0.116}
\providecommand{\combinedNResistant}{}
\renewcommand{\combinedNResistant}{125}
\providecommand{\combinedNPairs}{}
\renewcommand{\combinedNPairs}{4000}
\providecommand{\combinedNFullCombos}{}
\renewcommand{\combinedNFullCombos}{146}
\providecommand{\combinedNCombos}{}
\renewcommand{\combinedNCombos}{182}
\providecommand{\combinedTrainComboASR}{}
\renewcommand{\combinedTrainComboASR}{0.944}
\providecommand{\combinedTrainBudgetRatio}{}
\renewcommand{\combinedTrainBudgetRatio}{0.454}
}{}
\IfFileExists{../generated/signal_disentanglement_metrics.tex}
  {% Auto-generated by signal_disentanglement: DO NOT EDIT MANUALLY
% Last run: 2026-07-30 09:36:02

\providecommand{\signalAuthModelAUCB}{}
\renewcommand{\signalAuthModelAUCB}{0.881}
\providecommand{\signalAuthModelAUCA}{}
\renewcommand{\signalAuthModelAUCA}{0.860}
\providecommand{\signalAuthModelAUCC}{}
\renewcommand{\signalAuthModelAUCC}{0.832}
\providecommand{\signalMiaAUCBvsC}{}
\renewcommand{\signalMiaAUCBvsC}{0.6750}
\providecommand{\signalMiaAUCAvsC}{}
\renewcommand{\signalMiaAUCAvsC}{0.6575}
\providecommand{\signalMiaAUCBvsCLira}{}
\renewcommand{\signalMiaAUCBvsCLira}{0.5900}
\providecommand{\signalMiaAUCAvsCLira}{}
\renewcommand{\signalMiaAUCAvsCLira}{0.5125}
\providecommand{\signalMiaAUCBvsCCold}{}
\renewcommand{\signalMiaAUCBvsCCold}{0.6700}
\providecommand{\signalMiaAUCAvsCCold}{}
\renewcommand{\signalMiaAUCAvsCCold}{0.6575}
}{}

% ── metadata ──────────────────────────────────────────────────────────────────
\title{Membership Inference and Evasion Attacks\\on Gait-Based Authentication}
\author{Matteo Campagnaro}
\date{\today}

% ── header / footer ───────────────────────────────────────────────────────────
\pagestyle{fancy}
\fancyhf{}
\fancyhead[L]{\leftmark}
\fancyhead[R]{\thepage}
\renewcommand{\headrulewidth}{0.4pt}

% ══════════════════════════════════════════════════════════════════════════════
\begin{document}
% ══════════════════════════════════════════════════════════════════════════════

% ── title page ────────────────────────────────────────────────────────────────
\begin{titlepage}
  \centering
  \vspace*{2cm}
  {\Large University of Padova\par}
  {\large Department of Information Engineering\par}
  \vspace{1.5cm}
  {\LARGE\bfseries Membership Inference and Evasion Attacks\\[6pt]
   on Gait-Based Authentication\par}
  \vspace{2cm}
  {\large Master's Thesis in Cybersecurity\par}
  \vspace{2cm}
  \begin{tabular}{ll}
    \textbf{Author:}     & Matteo Campagnaro \\[6pt]
    \textbf{Supervisor:} & Prof.\ Simone Milani \\
  \end{tabular}
  \vfill
  {\large Academic Year 2025/2026\par}
\end{titlepage}

\pagenumbering{roman}

% ── abstract ──────────────────────────────────────────────────────────────────
\begin{abstract}
Gait-based authentication uses inertial sensor recordings collected during
ordinary device use to verify a user's identity, relying on the stability and
person-specificity of natural walking patterns. Deep learning models trained for
this task fit their training subjects closely enough to produce accurate
authentication decisions, but that same fitting introduces a privacy
vulnerability: the model's output behaviour may reveal which subjects were
present in the training set.

This thesis investigates membership inference attacks (MIA) on a gait
authentication pipeline consisting of a CNN encoder trained for subject
identification and a pairwise LSTM authentication head
that scores whether two gait windows come from the same person. Four attack variants are
evaluated across three inertial datasets (whuGAIT, UCI-HAR, WISDM) and a
cross-dataset configuration: a simple score-based baseline, a grey-box
likelihood ratio attack initialised from the target model weights, a warm-start
black-box variant, and a cold-start black-box variant. To determine how much of
the observed membership signal originates in the CNN encoder versus the LSTM
authentication head, a signal disentanglement experiment assigns subjects to
three non-overlapping groups so that the two components are trained on disjoint
populations.

Additionally, a projected gradient descent (PGD) evasion attack is evaluated on all
configurations, measuring the minimum perturbation budget required for a
successful impersonation relative to natural intra-subject gait variation.

Results show that a membership signal above chance is detectable across
configurations, with the authenticator and CNN encoder contributing comparable
amounts to the signal. Cold-start LiRA is the most consistent attack variant
across datasets. Evasion perturbations that succeed tend to stay within the
natural variation budget, raising practical concerns about the robustness of the
system.

\bigskip
\noindent\textit{Keywords:} gait authentication, membership inference attack,
LiRA, adversarial evasion, inertial sensors, privacy.
\end{abstract}

% ── front matter ──────────────────────────────────────────────────────────────
\tableofcontents
\listoffigures
\listoftables

\cleardoublepage
\pagenumbering{arabic}

% ── chapters ──────────────────────────────────────────────────────────────────
\chapter{Introduction}
\label{ch:introduction}

\section{Motivation}
\label{sec:intro_motivation}

The way a person walks is shaped by a combination of skeletal geometry, muscle coordination, and habitual movement patterns. These characteristics are person-specific enough to distinguish individuals reliably over short timescales, which makes gait a natural candidate for continuous biometric authentication: a system can verify a user's identity from inertial sensor readings collected during normal device use, without requiring any explicit interaction. Smartphones already carry the accelerometers and gyroscopes needed to collect gait signals, so the infrastructure cost of deployment is low.

Gait is, however, not static. Ageing, injury, weight change, and health conditions can alter a person's movement signature over months or years. This creates a practical problem for authentication systems that learn a fixed representation of a person's gait: the representation becomes stale, and the system must be updated without requiring the enrolled user to repeat the full enrolment procedure. This consideration motivates the two-component architecture used. The CNN encoder is trained once to learn a general-purpose gait representation from many subjects. The authentication head is trained separately to compare pairs of CNN-encoded windows and decide whether they come from the same person. Because authentication is framed as a pairwise distance comparison rather than a direct classification into a fixed set of enrolled identities, recalibration requires minimal retraining and fine-tuning of the authentication head on fresh reference windows from the same person, without touching the encoder. Adding a new user similarly requires no retraining of the encoder: only a set of reference windows is needed to enrol.

Any machine learning model trained on sensitive personal data carries a general privacy risk: models tend to fit their training examples more closely than held-out examples, and an attacker who can query the model may be able to exploit that asymmetry to infer whether a particular individual was in the training set. This is the membership inference problem~\cite{shokri2017membership}, and it is not specific to gait authentication. What is not known is how much membership signal is present in this particular architecture, under realistic threat models, and across different datasets and pipeline configurations. A person's gait encodes information about their age, health, and physical state, so the question has practical implications beyond the technical evaluation.

The layered architecture raises a further question that is specific to the two-component pipelines. Both the CNN encoder and the LSTM authentication head are trained on the same subjects in the standard configuration, so any memorisation is distributed across both components. It is therefore not clear how much of the observed membership signal, if any, comes from the encoder versus the authentication head. The two components have different architectures, different training objectives, and different amounts of exposure to each subject's data, so their individual contributions may differ substantially and may call for different mitigations. Separating them requires an experimental design in which the two components are trained on non-overlapping sets of subjects, which the signal disentanglement experiment in this thesis provides.

Beyond membership inference, a deployed authentication system faces a second question: how robust is its decision boundary? A small perturbation added to an impostor's probe signal, optimised via projected gradient descent (PGD), may be sufficient to push the model's output across the acceptance threshold. Measuring how large that perturbation needs to be, relative to the natural trial-to-trial variability of the signal, reveals how close the boundary sits to naturally occurring gait. Whether such perturbations are practically achievable, and whether they leave detectable spectral artefacts, has received less attention for gait-based systems than for image-domain authentication.

\section{Problem Statement}
\label{sec:intro_problem}

This thesis investigates two families of attacks on a gait-based authentication system built on the CNN architecture from Zou et al.~\cite{zou2020deep}.

The first family is membership inference. Given black-box or grey-box access to a trained authentication model, the attacker aims to determine whether a specific subject's gait data was part of the training set. The attack surface is the model's per-pair output score: the probability that two gait windows come from different people. A subject who was a training member may produce a systematically different score distribution compared to a held-out subject, and this difference constitutes the membership signal. The thesis evaluates several attack variants that differ in how much information the attacker has about the target model: a simple score-based baseline, a likelihood ratio attack using shadow models initialised from the target weights (grey-box), a warm-start black-box variant that uses a separately trained surrogate for initialisation, and a cold-start black-box variant that starts from random weights. The evaluation spans four dataset configurations: a within-dataset run on whuGAIT, within-dataset runs on UCI-HAR and WISDM, and a cross-dataset run in which a whuGAIT-pretrained encoder is paired with an authenticator retrained on subjects from all three datasets.

The second problem, addressed by a dedicated experiment, is the source of the membership signal. Since both the CNN encoder and the LSTM authentication head are trained on the same subjects in the standard pipeline, it is not possible to attribute the observed signal to one component or the other. To separate the two contributions, the thesis introduces a three-group experimental design in which subjects are partitioned so that the CNN sees one group, the authenticator sees a different group, and a third group serves as a held-out baseline seen by neither. Scoring all three groups with a single trained authenticator allows the per-group attack signal to be compared directly.

The third problem is evasion. For each dataset configuration, PGD is used to search for minimum-norm perturbations that cause the authentication model to accept an impostor's gait signal as genuine. The budget required for a successful attack is measured relative to the natural intra-subject variation in gait, providing a calibrated scale for the severity of the vulnerability.

\section{Contributions}
\label{sec:intro_contributions}

The main contributions of this thesis are as follows.

\begin{itemize}

  \item \textbf{Membership inference evaluation across multiple datasets and threat models.}
  The simple-delta attack and three LiRA variants (grey-box, warm-start black-box, cold-start black-box) are evaluated on four pipeline configurations: whuGAIT, UCI-HAR, WISDM, and a combined cross-dataset run. Cold-start consistently outperforms warm-start across all datasets, because random initialisation gives each shadow model an independent starting point and produces more diverse OUT distributions for the LiRA test. The delta variant comparison (raw, logit-first, mean-first) is therefore restricted to cold-start. The evaluation also identifies a grey-box signal inversion on whuGAIT whose cause is not fully established, and confirms cold-start LiRA as the least affected variant.

  \item \textbf{CNN vs authenticator memorisation signal disentanglement.}
  A three-group subject split is proposed in which the CNN encoder and the LSTM authenticator are trained on non-overlapping subjects within the whuGAIT dataset. Scoring all groups with a common authenticator allows the CNN's memorisation contribution to be separated from the LSTM's contribution. The experiment shows that both components produce signal above the held-out baseline, with the authenticator's contribution slightly exceeding the CNN's for simple-delta and cold-start LiRA, while the gap between the two narrows to near zero for the grey-box attack.

  \item \textbf{Evasion attack evaluation.}
  A PGD evasion attack is conducted on all four dataset configurations. The minimum perturbation budget required for a successful attack is compared against the median intra-subject gait variation, providing a naturalness baseline. Spectral analysis of the successful perturbations is included for the whuGAIT configuration.

\end{itemize}

\section{Thesis Outline}
\label{sec:intro_outline}

Chapter~\ref{ch:background} introduces the technical background: the properties of gait signals that make them suitable for authentication, the architecture of the CNN and LSTM components used in this work, and the formal framework for membership inference attacks including the LiRA construction.

Chapter~\ref{ch:related} surveys prior work on gait-based authentication, membership inference on biometric and machine learning systems, and adversarial attacks on authentication models, and positions this thesis relative to that literature.

Chapter~\ref{ch:system} describes the three datasets used in the evaluation, the subject split and pair construction procedure, the CNN encoder and authentication model architectures, and the four experimental configurations.

Chapter~\ref{ch:mia} presents the membership inference experiments: the simple-delta baseline, the LiRA shadow model attack in its three threat model variants, the delta formulation comparison, and a discussion of the grey-box signal inversion observed on whuGAIT. The chapter also presents the signal disentanglement experiment (Section~\ref{sec:mia_signal}): the three-group design, the authentication model performance per group, and the per-group membership signal comparison across attack variants.

Chapter~\ref{ch:evasion} presents the evasion attack: the PGD formulation, the epsilon budget estimation procedure, and the attack success rate and budget ratio results across all configurations, with a spectral analysis of the perturbations.

Chapter~\ref{ch:conclusions} summarises the findings, discusses the limitations of the experimental design, and identifies directions for future work.

\chapter{Background}
\label{ch:background}

This chapter introduces the concepts that will lay the foundation for the rest of the thesis. Section~\ref{sec:bg_gait} describes gait signals and how they are
captured with inertial sensors. Section~\ref{sec:bg_auth} explains the
distinction between identification and authentication and how deep learning
architectures are adapted for the authentication task.
Section~\ref{sec:bg_privacy} covers the connection between model memorisation
and privacy risk. Section~\ref{sec:bg_mia} formalises membership inference
attacks and describes the likelihood ratio formulation used in this work.

% ─────────────────────────────────────────────────────────────────────────────
\section{Gait Recognition}
\label{sec:bg_gait}

Gait recognition is the problem of identifying or verifying a person from the
way they walk. It can be approached from video, using the silhouette or
skeleton of a walking figure, or from inertial sensors attached to the body or
carried in a pocket. Sensor-based approaches are the more practical choice for
mobile authentication because the hardware is already present in any smartphone,
the measurement is passive and continuous, and no dedicated
infrastructures are required.

\subsection{Inertial Measurement Units and Gait Signals}
\label{ssec:bg_imu}

A modern smartphone contains at least a three-axis accelerometer and a
three-axis gyroscope, referred to as an inertial measurement unit
(IMU). The accelerometer measures specific force along each axis: when the
phone is stationary on a table it reads approximately $9.8\,\text{m/s}^2$
along the gravity axis, and when carried during walking it reads the
superposition of gravitational acceleration and the dynamic accelerations
produced by each step. The gyroscope measures angular velocity, which captures
the rotational motion of the device as the carrier shifts weight and swings
their arms.

Both sensors produce a continuous time series sampled at a fixed rate, typically
between 50\,Hz and 200\,Hz in consumer devices. The raw signal has a periodic
structure: each step produces a roughly repeatable pattern of peaks and troughs
whose shape, amplitude, and timing depend on the individual's body size, mass
distribution, stride length, and movement habits. This inter-subject variability
is the source of the biometric signal; the intra-subject stability over sessions
is what makes authentication feasible.

The phone's orientation and position relative to the body is not fixed, which is a practical complication. A phone carried in a trouser pocket sits differently from one held in the hand, and both differ from one lying in a bag. Several approaches address this: some datasets are collected with the device in a standardised position, some preprocessing pipelines project the signal onto a gravity-aligned frame, and some models are trained to be invariant to orientation through data augmentation. The datasets used in this thesis take different positions on this question, which is discussed in Chapter~\ref{ch:system}.

\subsection{Window Segmentation and Feature Representation}
\label{ssec:bg_windows}

Because gait is a periodic signal, the natural unit of analysis is a window
containing one or more complete gait cycles. A gait cycle is the interval
between two successive heel strikes of the same foot, lasting roughly one to
two seconds at a normal walking pace. Practical systems rarely segment by gait
cycle explicitly; instead they use fixed-length overlapping windows, which
avoids the need for reliable step detection and makes the windowing procedure
deterministic.

A window of length $T$ samples from $C$ sensor channels produces a matrix
$\mathbf{X} \in \mathbb{R}^{C \times T}$, which can be treated as a
multi-channel time series and fed directly into a convolutional or recurrent
network. Earlier approaches extracted handcrafted features from each window,
such as mean, variance, zero-crossing rate, and frequency-domain coefficients,
and then trained a classifier on the resulting feature vectors. Deep learning
pipelines learns a mapping from the raw window to a compact representation, with the representation shaped by the training objective.

The window length involves a trade-off. Longer windows capture more gait cycles and therefore more person-specific structure, improving discriminability between subjects. Shorter windows are more likely to be collected in practice, because any break in walking resets the signal and an authentication decision may be needed before a full long window is available. The datasets examined in this work use windows ranging from around two to five seconds, reflecting different positions on this trade-off.

% ─────────────────────────────────────────────────────────────────────────────
\section{Gait-Based Authentication}
\label{sec:bg_auth}

\subsection{Identification vs Authentication}
\label{ssec:bg_id_vs_auth}

Two related but distinct tasks arise in biometric systems: identification and
authentication. Identification is a one-to-many problem: given a measurement, determine which of $N$ known individuals it belongs to. The output is a class label, and the difficulty of the problem grows with $N$. This is the task used to train the CNN encoder in this work; training as a classifier with a softmax output over $N$ subject classes provides a supervision signal that encourages the learned representations to be discriminative across subjects.

Authentication is a one-to-one verification problem: given a measurement and a claimed identity, decide whether it matches the claim. The output is a binary accept or reject decision, and the difficulty does not depend on the total number of enrolled subjects in the same way. The device has one registered owner, and the question is whether the current user is that owner.

The two tasks use different metrics. Identification is evaluated by classification accuracy: the fraction of windows for which the predicted subject matches the true label. Authentication is evaluated by the false acceptance rate (FAR) and false rejection rate (FRR); the area under the (FAR-FRR) ROC curve (AUC) summarises overall separability and is the primary metric used throughout this thesis.

\subsection{Siamese and Contrastive Architectures}
\label{ssec:bg_siamese}

A natural architecture for authentication is a Siamese network~\cite{bromley1994signature}:
two copies of the same feature extractor, sharing weights, process the query
and the reference signal separately, and a comparison module produces a
similarity or distance score from the two resulting embeddings. If the score
exceeds a threshold, the query is accepted.

The key property of this architecture is that it does not require retraining
when a new subject is enrolled. The feature extractor is trained to produce
embeddings where same-subject pairs are close and different-subject pairs are
far apart, and any new subject can be enrolled simply by storing their reference
embedding. This is in contrast to a closed-set classifier, which requires
retraining to add a new class.

Training a Siamese network requires pairs and a loss that shapes the embedding
space. Common choices are contrastive loss~\cite{hadsell2006dimensionality} and
triplet loss~\cite{schroff2015facenet}; this thesis uses a simpler alternative
that treats each pair as a binary classification problem, with the model trained
with cross-entropy loss on the pair label and the output score interpreted as
$P(\text{different person} \mid \mathbf{x}_1, \mathbf{x}_2) \in [0, 1]$.

The system evaluated in this thesis uses this last approach. The CNN architecture
from Zou et al.~\cite{zou2020deep}, trained on the whuGAIT dataset, serves as
the shared feature extractor. The CNN is first trained as a subject identifier
on a set of training subjects. The authentication head, an LSTM followed by a
fully connected layer, is then trained on pairs of CNN-encoded windows with
cross-entropy loss. The score used for authentication decisions is
$P(\text{different person} \mid \mathbf{x}_1, \mathbf{x}_2)$. The CNN encoder
is kept frozen during authentication training, which keeps the memorisation
contributions of the two components separated and is discussed further in
Chapter~\ref{ch:signal}.

% ─────────────────────────────────────────────────────────────────────────────
\section{Privacy in Machine Learning}
\label{sec:bg_privacy}

\subsection{Memorisation and Generalisation}
\label{ssec:bg_memorisation}

Memorisation is distinct from overfitting: a model memorises a training example
if its output on that input depends on the example's presence in the training
set, regardless of whether the model generalises well on average.
Feldman~\cite{feldman2020does} showed that some memorisation of long-tail
examples is in fact necessary for optimal generalisation. The privacy implication
is that memorisation creates an exploitable asymmetry: if the model's behaviour
differs between training and non-training inputs, an attacker who can query the
model may detect that difference.

In gait authentication this asymmetry can manifest at multiple levels: individual pair scores, per-window confidence, or aggregated per-subject statistics. This thesis focuses on the subject-level signal, using the gap between training and non-training subjects as the membership signal and establishing a baseline attack on an unprotected system, without evaluating privacy-hardened mitigations.

% ─────────────────────────────────────────────────────────────────────────────
\section{Membership Inference Attacks}
\label{sec:bg_mia}

\subsection{Threat Model}
\label{ssec:bg_threat}

A membership inference attack targets a trained model $f_\theta$ and asks,
for a given data point $x$ and label $y$: was $(x, y)$ in the training set
$\mathcal{D}_\text{train}$ used to produce $\theta$? The attacker's goal is to
build a binary classifier $\mathcal{A}$ such that $\mathcal{A}(x, y, f_\theta)
= 1$ indicates membership and $0$ indicates non-membership.

The attacker's capabilities vary along two dimensions. The first is access to
the model: a white-box attacker can read the model weights and gradients
directly; a grey-box attacker can query the model and observe its output
probabilities but knows the architecture and possibly the training procedure;
a black-box attacker can only query the model and observe its outputs, with no
knowledge of the weights or training procedure. The second dimension is
knowledge of the training data distribution: the attacker may have access to
data from the same distribution (shadow data) without knowing the exact training
set.

For gait authentication, the relevant queries are pairs of gait windows. The
model's output for a pair $(\mathbf{x}_1, \mathbf{x}_2)$ is a scalar
$p = P(\text{different person} \mid \mathbf{x}_1, \mathbf{x}_2) \in [0, 1]$.
The attacker has access to a set of same-person and different-person pairs for
a target subject $s$ and observes the model's score distribution over those
pairs. The membership question is then: was subject $s$'s data in the
authentication model's training set?

This thesis evaluates attacks under three threat models. The grey-box attacker
knows the model architecture and initialises shadow models from a copy of the
target weights, which corresponds to an attacker who has obtained the model
checkpoint through leakage or reverse engineering. The warm-start black-box
attacker trains a surrogate model independently and uses its weights for shadow
model initialisation, corresponding to an attacker with knowledge of the
architecture and training procedure but not the weights. The cold-start
black-box attacker uses random initialisation, corresponding to the most
restrictive assumption in which the attacker knows only the architecture.

\subsection{Shadow Model Approaches}
\label{ssec:bg_shadow}

The shadow model technique, introduced by Shokri et al.~\cite{shokri2017membership},
is the foundation of most subsequent MIA work. The idea is to train several
models on datasets drawn from the same distribution as the target model's
training set, with known membership labels: for each shadow model, the attacker
knows which examples were in its training set and which were held out. The
shadow models produce output distributions for both member and non-member
examples, and an attack classifier is trained on those distributions to
distinguish the two cases. At test time the same attack classifier is applied
to the target model's output on the query example.

The effectiveness of this approach depends on the shadow models approximating
the target model's behaviour on members and non-members. If the target model
was trained differently from the shadow models, or on a substantially different
data distribution, the shadow distributions will not match and the attack will
underperform.

A cleaner formulation replaces the learned attack classifier with a direct
statistical test. For a given example, the attacker collects scores from many
shadow models: scores from shadow models that included the example in their
training set (IN scores) and scores from shadow models that excluded it (OUT
scores). These two sets of scores approximate the conditional distributions
$P(\text{score} \mid \text{member})$ and $P(\text{score} \mid
\text{non-member})$. The membership decision is made by comparing the target
model's score on the example against these two distributions.

\subsection{Likelihood Ratio Attacks (LiRA)}
\label{ssec:bg_lira}

Carlini et al.~\cite{carlini2022membership} formalised this statistical
formulation as a likelihood ratio attack (LiRA). For a target example $z =
(x, y)$ and a target model $f_\theta$, let $\phi(z, \theta)$ be the
per-example score (e.g. the model's loss or confidence on $z$). The attacker
trains $K$ shadow models $\theta_1^{\text{in}}, \ldots, \theta_K^{\text{in}}$
with $z$ included in each training set and $K$ shadow models
$\theta_1^{\text{out}}, \ldots, \theta_K^{\text{out}}$ with $z$ excluded. The
observed scores
\begin{align}
  \phi^{\text{in}}_k  &= \phi(z, \theta_k^{\text{in}}),  \quad k = 1, \ldots, K, \\
  \phi^{\text{out}}_k &= \phi(z, \theta_k^{\text{out}}), \quad k = 1, \ldots, K,
\end{align}
are used to fit Gaussian distributions $\mathcal{N}(\mu^{\text{in}},
\sigma^{\text{in}})$ and $\mathcal{N}(\mu^{\text{out}}, \sigma^{\text{out}})$
by maximum likelihood. The membership score for $z$ under the target model is
the log-likelihood ratio:
\begin{equation}
  \Lambda(z) = \log \frac{p(\phi(z, \theta) \mid \mathcal{N}(\mu^{\text{in}},
  \sigma^{\text{in}}))}{p(\phi(z, \theta) \mid \mathcal{N}(\mu^{\text{out}},
  \sigma^{\text{out}}))}.
  \label{eq:lira}
\end{equation}
A positive value of $\Lambda(z)$ indicates that the target score is more
consistent with the in-distribution than the out-distribution, suggesting
membership.

LiRA is statistically principled in the sense that, if the Gaussian assumption
holds and the shadow models are drawn from the same distribution as the target,
the likelihood ratio is the optimal test for membership at any false positive
rate~\cite{carlini2022membership}. In practice, the Gaussian fit is
approximate, the shadow model distribution may differ from the target, and $K$
is finite, all of which introduce error.

\noindent\textbf{Online vs offline LiRA.}
The formulation above is the \emph{online} variant: for each target example,
half of the $K$ shadow models are trained with it included (IN) and half with
it excluded (OUT), producing two empirical distributions. This is expensive
because a separate set of shadow models must be trained for each target example.

The \emph{offline} variant avoids this cost by training all $K$ shadow models
without the target example, producing only the OUT distribution
$\mathcal{N}(\mu^{\text{out}}, \sigma^{\text{out}})$. The membership score
then becomes a one-sided test: how far above the OUT baseline does the target
model's score fall?
\begin{equation}
  \Lambda_{\text{offline}}(z)
  = \Phi\!\left(\phi(z, \theta);\, \mu^{\text{out}},\, \sigma^{\text{out}}\right),
  \label{eq:lira_offline}
\end{equation}
where $\Phi(\,\cdot\,; \mu, \sigma)$ is the Gaussian CDF. A score close to 1
indicates that the target model's value on $z$ is unusually high relative to
what shadow models produce when $z$ is absent from their training data,
suggesting membership. A score near 0.5 is consistent with the OUT distribution
and suggests non-membership.

The offline variant is used throughout this thesis. Its practical advantage is
that the same $K$ shadow models can be reused for every target subject, because
none of them ever includes any target subject's data. This is a natural fit for
the subject-level attack: the shadow models are trained on random subsets of
available pair data, and each target subject's pairs are simply held out of all
shadows. The statistical cost relative to the online variant is small when the
target example has limited influence on each shadow, which holds when the
training set is large relative to a single subject's contribution.

For gait authentication, the per-example score $\phi$ is adapted to the
pairwise nature of the model. Rather than a single loss value for a single
example, the relevant quantity for a subject $s$ is a per-subject delta:
\begin{equation}
  \delta(s) = \bar{P}_{\text{imp}}(s) - \bar{P}_{\text{gen}}(s),
  \label{eq:bg_delta}
\end{equation}
where $\bar{P}_{\text{imp}}(s)$ is the mean model output $P(\text{different
person})$ over impostor pairs involving $s$, and $\bar{P}_{\text{gen}}(s)$ is
the mean over genuine pairs. A training subject whose gait representation has
been fitted tightly by the model will tend to produce a larger gap between
impostor and genuine scores, resulting in a higher $\delta(s)$ compared to a
held-out subject. The shadow models are trained with and without each target
subject's pairs, and the resulting distributions of $\delta$ values are used in
place of the standard loss distributions in Equation~\eqref{eq:lira}.

\chapter{Related Work}
\label{ch:related}

% ─────────────────────────────────────────────────────────────────────────────
\section{Gait-Based Biometric Authentication}
\label{sec:rel_gait_auth}

\subsection{Early Approaches and Feature Engineering}
\label{ssec:rel_gait_early}

Interest in gait as a biometric signal dates to the early 2000s, when
Gafurov et al.~\cite{gafurov2007gait} demonstrated that accelerometer recordings
from waist-mounted devices could distinguish individuals at rates well above
chance using handcrafted cycle-level features. These features proved brittle
when device placement or walking speed varied. The UCI-HAR~\cite{anguita2013human}
and WISDM~\cite{kwapisz2011activity} datasets were originally collected for
activity recognition but became common benchmarks for gait-related tasks due to
their multi-subject, naturalistic recordings.

\subsection{Deep Learning Approaches}
\label{ssec:rel_gait_deep}

Zou et al.~\cite{zou2020deep} showed that a CNN encoder trained directly on raw sensor windows captures person-specific temporal patterns with sufficient discriminability for both identification and pairwise authentication, without requiring explicit step detection. They combine the CNN encoder with an LSTM authentication head trained on pairs of CNN-encoded windows, decoupling representation learning from pairwise comparison. This two-stage architecture is the basis of the system evaluated in this thesis.

\subsection{Privacy in Gait Systems}
\label{ssec:rel_gait_privacy}

Privacy analysis of trained gait models has received considerably less attention
than recognition accuracy. Malekzadeh et al.~\cite{malekzadeh2018mobile} address
gait anonymisation, a different threat model aimed at protecting raw signals
before transmission rather than limiting what a trained model leaks.
Milani~\cite{milani2024gait} analyses identity inference vulnerability in a gait
identification pipeline. The present thesis extends that analysis to an
authentication pipeline and evaluates it across multiple datasets.

% ─────────────────────────────────────────────────────────────────────────────
\section{Membership Inference on Machine Learning Models}
\label{sec:rel_mia_biometric}

\subsection{Foundational Work}
\label{ssec:rel_mia_foundations}

Shokri et al.~\cite{shokri2017membership} formalised membership inference by
training shadow models with known membership labels and using their outputs to
build an attack classifier. Yeom et al.~\cite{yeom2018privacy} grounded the
attack theoretically: the training-test loss gap upper-bounds the attack
advantage, directly connecting memorisation to vulnerability. Salem
et al.~\cite{salem2019ml} showed that a single shadow model and no knowledge
of the target architecture are often sufficient, lowering the bar for practical
attacks. Carlini et al.~\cite{carlini2022membership} unified these ideas into a
likelihood ratio framework (LiRA) that is statistically optimal at any false
positive rate and substantially outperforms earlier methods, particularly on
models with small generalisation gaps; the offline variant is the formulation
adopted in this thesis and is described in detail in
Section~\ref{ssec:bg_lira}.

\subsection{MIA on Biometric Models}
\label{ssec:rel_mia_biometric_detail}

Standard MIA pipelines operate on the softmax output of a classifier, using
the per-class confidence vector as the membership signal. Authentication models
produce no such output: the model scores pairs of windows rather than assigning
class probabilities to individual inputs, so the standard attack surface does
not exist. He and Zhang~\cite{he2021quantifying} studied this problem in the
context of contrastive learning, where models similarly produce embeddings
rather than class probabilities. They adapted existing MIA methods to operate
on the representation space and found that contrastive models are less
vulnerable to membership inference than supervised classifiers, attributing
this to their reduced tendency to overfit individual training examples. This
thesis faces the same structural obstacle and derives an alternative membership
signal from the distribution of pairwise authentication scores, as described
in Chapter~\ref{ch:mia}.

% ─────────────────────────────────────────────────────────────────────────────
\section{Adversarial Attacks on Authentication}
\label{sec:rel_evasion}

\subsection{Adversarial Examples}
\label{ssec:rel_evasion_foundations}

Szegedy et al.~\cite{szegedy2014intriguing} identified that small input
perturbations can reliably change a classifier's output. Goodfellow
et al.~\cite{goodfellow2015explaining} proposed FGSM as a single-step gradient
attack; Madry et al.~\cite{madry2018towards} showed that iterative PGD with
random restart finds stronger examples and serves as a principled robustness
benchmark. The evasion attack in this thesis uses PGD adapted to the pairwise
authentication setting: the goal is an impostor gait window that the model
accepts as the enrolled target.

\subsection{Adversarial Attacks on Biometric Systems}
\label{ssec:rel_evasion_biometric}

Sharif et al.~\cite{sharif2016accessorize} demonstrated physical adversarial
attacks on face recognition via printed eyeglass frames. Gait authentication
presents a different attack surface: the signal is a sensor recording, so
perturbations are applied digitally rather than physically. This thesis
evaluates such digital attacks and assesses whether the required perturbation
budget is consistent with natural within-subject gait variability.

% ─────────────────────────────────────────────────────────────────────────────
\section{Positioning of This Work}
\label{sec:rel_positioning}

This thesis applies membership inference to a pairwise authentication model,
where the standard softmax-based attack surface does not exist and the
membership signal must be derived from the distribution of pair scores. The
evaluation spans multiple datasets under three controlled threat models.
A signal disentanglement experiment isolates the privacy contribution of the
CNN encoder from that of the LSTM authentication head, treating the two
components separately by construction. The evasion evaluation complements
attack success rate with a check against natural within-subject gait
variability, giving a concrete sense of whether successful perturbations
are practically realisable.

\chapter{System Model and Datasets}
\label{ch:system}

The pipeline has two sequential learning stages: a CNN encoder trained to classify subjects by their gait, and an LSTM authentication head that scores pairs of gait windows. The CNN is trained first and then frozen; only the LSTM is updated during authentication training. Three inertial datasets are used, each with its own partition of subjects into members and non-members, and pair inventories are constructed to keep the two groups strictly separate.

Because both components are trained on the same subjects in the standard configurations, it is not possible to attribute the observed membership signal to one or the other. To address this, a dedicated experiment trains the CNN and the LSTM on entirely disjoint sets of subjects within whuGAIT, so that the privacy leakage of each component can be measured independently. This experiment is a standalone setup described in Section~\ref{ssec:sys_signal} and analysed in Section~\ref{sec:mia_signal}.

% ─────────────────────────────────────────────────────────────────────────────
\section{Datasets}
\label{sec:sys_datasets}

Three publicly available inertial datasets are used: whuGAIT, UCI-HAR, and WISDM. All three provide accelerometer and gyroscope recordings of walking collected with smartphones or waist-mounted devices. They differ in sampling rate, collection protocol, number of subjects, and the activities recorded. Table~\ref{tab:datasets} summarises their main properties.

\begin{table}[htbp]
\centering
\caption{Dataset summary.}
\label{tab:datasets}
\begin{tabular}{lrrrrl}
\toprule
Dataset  & Subjects & Sampling & Window    & Window    & Segmentation \\
         &          & rate     & samples   & duration  & \\
\midrule
whuGAIT  & 118      & 50\,Hz   & 128       & 2.56\,s   & Step-based \\
UCI-HAR  & 30       & 50\,Hz   & 128       & 2.56\,s   & Pre-segmented \\
WISDM    & 51       & 20\,Hz   & 128       & 6.40\,s   & Sliding window (50\%) \\
\bottomrule
\end{tabular}
\end{table}

\subsection{whuGAIT}
\label{ssec:sys_whugait}

The whuGAIT dataset~\cite{zou2020deep} contains recordings from 118 subjects collected in the wild: participants carried their smartphones in an unrestricted manner while walking along arbitrary paths indoors and outdoors. The device's three-axis accelerometer and three-axis gyroscope are sampled at 50\,Hz, giving six channels per timestep. Windows of 128 samples (2.56 seconds) are extracted by detecting individual steps in the raw signal and aligning windows to the detected step boundaries, which makes each window more likely to contain a complete gait cycle than a purely time-based approach.

The dataset is distributed with a fixed pair inventory (Dataset \#5 in the original paper) in which same-person and different-person pairs are pre-formed from the 128-sample windows. Each entry in the inventory consists of two windows and a binary label indicating whether the two windows belong to the same person or to different people. This fixed structure is used directly in the authentication training and attack evaluation, rather than constructing pairs from scratch.

\subsection{UCI-HAR}
\label{ssec:sys_ucihar}

The UCI Human Activity Recognition dataset~\cite{anguita2013human} contains recordings from 30 subjects performing six labelled activities: walking, walking upstairs, walking downstairs, sitting, standing, and lying. The device is carried on the waist at 50\,Hz. Windows of 128 samples (2.56 seconds) are pre-segmented and provided with activity labels. Only the walking activity windows are retained for the gait authentication task, as the other activities would introduce activity-level variation that is not a property of an individual's gait. The dataset provides a predefined subject-level split into 21 training subjects and 9 test subjects, which is used directly as the member/non-member partition.

\subsection{WISDM}
\label{ssec:sys_wisdm}

The WISDM dataset~\cite{kwapisz2011activity} contains accelerometer and gyroscope recordings from 51 subjects using a smartphone placed in a trouser pocket, sampled at approximately 20\,Hz. Six activities are labelled; only the walking activity is retained. Because the recordings are continuous rather than pre-segmented, windows of 128 samples (6.4 seconds at 20\,Hz) are extracted with a sliding window of stride 64, giving 50\% overlap between consecutive windows. The lower sampling rate means the 128-sample windows cover a longer time interval than in the 50\,Hz datasets, capturing more gait cycles per window.

% ─────────────────────────────────────────────────────────────────────────────
\section{Subject Split and Pair Construction}
\label{sec:sys_split}

\subsection{Member and Non-Member Assignment}
\label{ssec:sys_membership}

For each dataset a fixed subject-level partition assigns subjects to one of two roles for the purposes of the MIA evaluation. \emph{Member} subjects are those whose data appears in the training set of at least one of the two pipeline stages: CNN identification training or authentication head training. \emph{Non-member} subjects are held out from training entirely and serve as the attack's negative class.

The sizes of these partitions reflect the total number of subjects in each dataset: for whuGAIT, 98 subjects are members and 20 are non-members; for UCI-HAR, 21 are members and 9 are non-members (following the predefined dataset split); for WISDM, 40 are members and 11 are non-members.

For the within-dataset configurations (whuGAIT, UCI-HAR, WISDM), the CNN encoder is trained from scratch on the member subjects of that dataset, and the authentication head is subsequently trained on the same member set. For the combined configuration, the CNN checkpoint trained on whuGAIT members is reused as a frozen encoder, and only the authentication head is trained on members pooled across all three datasets. This is discussed further in Section~\ref{sec:sys_configs}.

\subsection{Authentication Pair Inventory}
\label{ssec:sys_pairs}

For the whuGAIT dataset the fixed pair inventory from the original paper (D5) is used directly. D5 is distributed as two separate directories: a training set of 66,542 pairs involving subjects 21--118 and a test set of 7,600 pairs involving subjects 1--20; the member assignment is defined to match D5's existing partition.

For UCI-HAR and WISDM no pre-built pair inventory exists, so one is constructed from the raw windows. The datasets ship with a predefined subject-level train/test split that is disjoint by construction, which is used as the member/non-member partition. Pairs are then sampled independently from the two window pools: member pairs are drawn from member windows only, and non-member pairs from non-member windows only, so the two groups never share windows. Each inventory is balanced between same-person and different-person pairs to avoid class imbalance in authentication training.

For the MIA attack, pair membership must be consistent with the subject membership label: a pair is a member pair if and only if both subjects in the pair are members. This \emph{member-aware} pair split ensures that the delta signal $\delta(s) = \bar{P}_\text{imp}(s) - \bar{P}_\text{gen}(s)$ is computed from pairs whose membership status is unambiguous.

% ─────────────────────────────────────────────────────────────────────────────
\section{CNN Encoder}
\label{sec:sys_cnn}

\subsection{Architecture}
\label{ssec:sys_cnn_arch}

The CNN encoder follows the architecture from Zou et al.~\cite{zou2020deep}. Figure~\ref{fig:cnn_arch} shows the full layer sequence with output shapes at each stage. The input window $(B, 6, 128)$ is reshaped to $(B, 1, 6, 128)$, treating the six sensor channels as spatial height and the 128 time steps as width of a single-channel image. The first three convolutional blocks operate along the time axis only, learning temporal patterns independently per channel. The fourth block uses a $6 \times 1$ kernel that collapses all channels at once, fusing the per-channel temporal patterns into a joint representation of shape $(B, 128, 1, 16)$.

\begin{figure}[htbp]
  \centering
  \includegraphics[width=\textwidth]{cnn_architecture.png}
  \caption{CNN encoder architecture (Zou et al.~\cite{zou2020deep}). Output
           shapes are shown below each block (batch dimension omitted). The
           dashed branch indicates the feature maps $(16, 128)$ passed to the
           LSTM authentication head; the main path continues to the FC layer
           used during identification pre-training.}
  \label{fig:cnn_arch}
\end{figure}

For identification pre-training the feature map is flattened to a 2048-dimensional vector and passed through a fully connected layer to produce logits over the $N$ training subjects. For the authentication head the flattening and FC layer are bypassed: the feature map is transposed to $(B, 16, 128)$ and fed directly to the LSTM.

\subsection{Training as Subject Identifier}
\label{ssec:sys_cnn_train}

For each within-dataset configuration the CNN is trained as a subject classifier on that dataset's member subjects, minimising cross-entropy loss on the per-window subject label. Training uses the Adam optimiser with a learning rate of $2.5 \times 10^{-3}$ and L2 weight decay of $1.5 \times 10^{-3}$. The CNN weights are fixed after identification training and are not updated during any subsequent stage.

% ─────────────────────────────────────────────────────────────────────────────
\section{Authentication Model}
\label{sec:sys_auth}

\subsection{LSTM Head and Similarity Score}
\label{ssec:sys_lstm}

The authentication model pairs the frozen CNN encoder with a trainable LSTM head, as shown in Figure~\ref{fig:auth_arch}. Both windows are encoded to $(16, 128)$ feature sequences, concatenated to $(32, 128)$, and processed by a 2-layer LSTM (hidden size 64). The last hidden state is projected through dropout and a fully connected layer to a score $P(\text{different person} \mid \mathbf{x}_1, \mathbf{x}_2) \in [0, 1]$; a pair is accepted as genuine if this score is below 0.5. The LSTM head has 83,074 trainable parameters.

\begin{figure}[htbp]
  \centering
  \includegraphics[width=\textwidth]{auth_architecture.png}
  \caption{AuthModel architecture. Both windows pass through the same frozen
           CNN encoder (weights shared). The resulting $(16,128)$ sequences
           are concatenated and fed to a 2-layer LSTM; the last hidden state
           is projected to a similarity score via dropout and a fully connected
           layer.}
  \label{fig:auth_arch}
\end{figure}

\subsection{Training Procedure}
\label{ssec:sys_auth_train}

The LSTM head is trained on the member pairs using the Adam optimiser with learning rate $2.5 \times 10^{-3}$, L2 weight decay $1.5 \times 10^{-3}$, batch size 512, and dropout rate 0.3. Training runs for 10 epochs. The CNN weights are frozen throughout; only the LSTM weights and the FC layer are updated. This keeps the two stages independent: the CNN is shaped by the identification objective, the LSTM by the pairwise comparison objective. The signal disentanglement experiment in Section~\ref{sec:mia_signal} exploits this separation to measure how much each component contributes to the membership signal.

Figure~\ref{fig:training_curves} shows the per-epoch train and test accuracy and loss for the authentication head across all four experimental configurations. The dotted vertical line marks the epoch selected by early stopping. whuGAIT trains for 2 epochs before the test accuracy stagnates; UCI-HAR and the combined run stop at epochs 2 and 8 respectively; WISDM trains for the full 10 epochs with the test accuracy still improving at the end. The growing train/test gap in UCI-HAR is consistent with the relatively small member pool (21 subjects), which gives the LSTM fewer pairs to generalise from.

\begin{figure}[htbp]
  \centering
  \includegraphics[width=\textwidth]{analysis/models_overfitting_curves.png}
  \caption{Per-epoch train (solid) and test (dashed) accuracy and cross-entropy
           loss for the authentication head across all four configurations. The
           dotted vertical line marks the selected checkpoint epoch.}
  \label{fig:training_curves}
\end{figure}

% ─────────────────────────────────────────────────────────────────────────────
\section{Experimental Configurations}
\label{sec:sys_configs}

Five configurations are evaluated. The first four differ in which dataset the CNN and authentication head are trained on and which subjects form the held-out test set. The fifth is the signal disentanglement configuration, in which the CNN and LSTM are trained on disjoint subject groups within whuGAIT to isolate each component's contribution to the membership signal (Section~\ref{sec:mia_signal}).

\subsection{Within-Dataset Runs}
\label{ssec:sys_within}

The three within-dataset configurations each use a single dataset for all stages of the pipeline.

\textbf{whuGAIT.} The CNN is trained for identification on the 98 member subjects. The authentication head is then trained on pairs drawn from the same 98 subjects. The 20 non-member subjects serve as the held-out class for both the authentication evaluation and the MIA attack.

\textbf{UCI-HAR.} The CNN is trained from scratch on the 21 UCI-HAR member subjects for identification, then frozen. The authentication head is trained on pairs from the same 21 subjects. The 9 held-out subjects are the non-member class.

\textbf{WISDM.} The CNN is trained from scratch on the 40 WISDM member subjects for identification, then frozen. The authentication head is trained on pairs from the same 40 subjects, with 11 subjects held out.

\subsection{Cross-Dataset (Combined) Run}
\label{ssec:sys_combined}

The combined configuration assembles a training set from all three datasets: the 98 whuGAIT members, the 21 UCI-HAR members, and the 40 WISDM members, giving 159 training subjects in total. The whuGAIT CNN encoder is used and kept frozen. The authentication head is trained on pairs from all 159 subjects simultaneously. The resulting training inventory contains 72,000 pairs: 60,000 from whuGAIT and 6,000 each from UCI-HAR and WISDM. The imbalance reflects the larger subject count of whuGAIT, since pairs are sampled proportionally per subject.

The held-out test set consists of the 9 UCI-HAR non-members and the 11 WISDM non-members (20 subjects total), identified by a numeric offset that encodes their source dataset. No whuGAIT non-members are included in the combined test set, because the whuGAIT CNN was pre-trained on the whuGAIT members and the MIA signal for whuGAIT subjects is analysed separately in the within-dataset run.

\subsection{Signal Disentanglement Run}
\label{ssec:sys_signal}

The signal disentanglement configuration is a fifth pipeline run on whuGAIT using 60 of the 118 available subjects. The experiment requires three non-overlapping groups of equal size to ensure that any difference in measured membership signal does not depend on the training set size: Group C is fixed at subjects 1--20, which is the existing D5 non-member partition, so Groups A and B are also set to 20 subjects each to ensure that the CNN and the LSTM are trained on the same amount of data. The remaining 58 whuGAIT subjects are not used in this configuration. The three groups are defined as follows, so that the CNN encoder and the LSTM authentication head are trained on entirely disjoint populations:

\begin{itemize}
  \item \textbf{Group A (subjects 21--40):} used only for CNN identification
    training. The authentication head never sees pairs from these subjects.
  \item \textbf{Group B (subjects 41--60):} used only for authentication head
    training. The CNN is not trained to classify these subjects.
  \item \textbf{Group C (subjects 1--20):} held out from both training stages
    and used as the non-member baseline for the MIA attack.
\end{itemize}

Because the CNN and LSTM are trained on different subjects, the membership signal produced by each component is anchored to a different population. Comparing the attack's ability to distinguish Group A from Group C isolates the CNN's contribution, while comparing Group B against Group C isolates the authenticator's contribution. The two comparisons are evaluated on the same model and the same held-out group, making them directly comparable. The results are analysed in Section~\ref{sec:mia_signal}.
\chapter{Membership Inference Attack}
\label{ch:mia}

The central question of this chapter is whether the authentication model's output behaviour reveals which subjects were in its training set. The attack score is built from the per-subject delta $\delta(s)$, which measures how much more discriminatively the model separates a subject's genuine and impostor pairs. A simple threshold on $\delta(s)$ already provides a baseline attack that requires no knowledge of the model's internals. To go further, LiRA casts the problem as a likelihood ratio test by training shadow models and comparing the target subject's score distribution against what would be expected if they were absent from training. Three shadow-model initialisation strategies are evaluated. The grey-box variant assumes access to the target checkpoint. The warm-start variant tests whether a shared surrogate initialisation can reduce per-shadow training cost without substantially sacrificing attack effectiveness. The cold-start variant uses independent random initialisation for each shadow and consistently outperforms the other two. Once cold-start is established as the most effective threat model, three delta formulations are compared on it. On whuGAIT, both grey-box and warm-start produce inverted signals whose cause is not fully established.

% ─────────────────────────────────────────────────────────────────────────────
\section{Attack Formulation}
\label{sec:mia_formulation}

Following the Identity Inference Attack (IIA) framing of Milani~\cite{milani2024gait}, the goal is to determine whether a target subject $s$ was enrolled in the model's training set, using only gait windows from $s$ that need not be the exact training samples~\cite{li2022userlevel}. The attack score is derived from the model's pairwise output distribution for $s$.

\subsection{Per-Subject Delta Score}
\label{ssec:mia_delta}

The authentication model takes a pair of gait windows as input and outputs $P(\text{different person} \mid \mathbf{x}_1, \mathbf{x}_2) \in [0, 1]$. A high output means the model considers the pair to come from two different people; a low output means it considers them genuine. For a subject $s$, two distributions of such scores can be observed:

\begin{itemize}
  \item \emph{Genuine pairs}: pairs in which both windows come from $s$.
    A well-fitted model should score these low (close to 0).
  \item \emph{Impostor pairs}: pairs in which one window comes from $s$ and
    one from a different subject.
    A well-fitted model should score these high (close to 1).
\end{itemize}

The per-subject delta is the mean gap between these two distributions (see Equation~\eqref{eq:bg_delta} in Section~\ref{ssec:bg_lira}):
\begin{equation*}
  \delta(s) = \bar{P}_{\text{imp}}(s) - \bar{P}_{\text{gen}}(s),
\end{equation*}
where $\bar{P}_{\text{imp}}(s)$ is the mean model output over impostor pairs involving $s$ and $\bar{P}_{\text{gen}}(s)$ is the mean over genuine pairs. Using raw $P(\text{different person})$ as the membership signal directly is not straightforward: the signal direction is opposite for the two pair types, since a well-fitted model drives genuine scores down and impostor scores up. A single pair score is therefore ambiguous without knowing the pair type, and using only one pair type discards half the available signal. The delta resolves this by combining both in a single consistent quantity. It also cancels subject-level baseline differences: some subjects have inherently more distinctive gait and produce high impostor scores regardless of training membership; subtracting $\bar{P}_{\text{gen}}$ normalises for this, so $\delta(s)$ measures how much more separated the model is for $s$ relative to that subject's own baseline.

\subsection{Member vs Non-Member Hypothesis}
\label{ssec:mia_hypothesis}

The membership hypothesis is that the model fits its training subjects more tightly than held-out subjects, producing systematically larger deltas for members. Concretely, if $\mathcal{M}$ and $\mathcal{N}$ denote the sets of member and non-member subjects respectively, the attack expects:
\begin{equation}
  \mathbb{E}_{s \in \mathcal{M}}[\delta(s)]
  > \mathbb{E}_{s \in \mathcal{N}}[\delta(s)].
\end{equation}
An attack that uses $\delta(s)$ directly or a function of it to classify $s$ as member or non-member is evaluated by the area under the ROC curve (AUC) over the binary classification of all subjects. AUC $= 0.5$ corresponds to random guessing; AUC $= 1.0$ corresponds to perfect separation. The TPR at 10\% FPR is reported alongside AUC to give a consistent low-FPR operating point that is comparable across all attack variants.

% ─────────────────────────────────────────────────────────────────────────────
\section{Simple-Delta Attack}
\label{sec:mia_simple}

\subsection{Results Across Datasets}
\label{ssec:mia_simple_results}

The simplest attack uses $\delta(s)$ directly as the membership score, with no shadow models. A subject $s$ is classified as a member if $\delta(s)$ exceeds a threshold. This attack requires only black-box query access to the model and knowledge of which windows belong to which subject. Members tend to produce higher deltas because the LSTM has been trained directly on their pairs, making it more discriminative for their gait patterns. Non-members have never been seen during training, so their genuine and impostor pairs are scored with less precision and the gap is smaller.

Table~\ref{tab:simple_delta} reports the simple-delta AUC and TPR at 10\% FPR for all four configurations. Because the non-member population ranges from 9 to 20 subjects, the ROC curve is coarse and the TPR at 10\% FPR value can only take a small number of discrete steps; the reported figure should be read as an approximate operating point rather than a precise measurement.

\begin{table}[htbp]
\centering
\caption{Simple-delta MIA results. TPR at 10\% FPR is coarse due to small non-member populations (9--20 subjects).}
\label{tab:simple_delta}
\begin{tabular}{lrr}
\toprule
Configuration & AUC & TPR at 10\% FPR \\
\midrule
whuGAIT  & \whuGAITSimplDAUC  & \whuGAITSimplDTPRten \\
UCI-HAR  & \uciharSimplDAUC   & \uciharSimplDTPRten  \\
WISDM    & \wisdmSimplDAUC    & \wisdmSimplDTPRten   \\
Combined & \combinedSimplDAUC & \combinedSimplDTPRten \\
\bottomrule
\end{tabular}
\end{table}

The simple-delta attack produces AUC above 0.5 in all four configurations, confirming that a membership signal is present in the authentication model across datasets. WISDM shows the strongest signal (\wisdmSimplDAUC), reflecting the fact that 128-sample windows at 20\,Hz span 6.4 seconds and therefore capture several complete gait cycles per window, giving the model more person-specific structure to memorise. The combined configuration shows the weakest signal (\combinedSimplDAUC), consistent with the cross-dataset distribution shift making it harder for the LSTM to fit any single subject's gait as tightly as in the within-dataset runs.

% ─────────────────────────────────────────────────────────────────────────────
\section{LiRA Shadow Model Attack}
\label{sec:mia_lira}

\subsection{Shadow Model Training}
\label{ssec:mia_shadow_train}

The LiRA attack requires training $K$ shadow LSTM models, each on a different subset of the training pairs. The shadow models approximate the distribution of model behaviours when a target subject is excluded from training (OUT distribution), which is then used to assess how unusual the target model's delta is for that subject.

Each shadow is trained with the member-aware split protocol. For each shadow $k$, a random 50\% of the known member subjects is drawn as the included set. Only pairs whose both windows come from subjects in the included set are used to train that shadow. This ensures that a target subject assigned as OUT for shadow $k$ does not appear in its training data even through the reference window of a cross-subject pair. The remaining 50\% of members constitutes the OUT set for that shadow, and the shadow's delta values for those OUT subjects form the empirical samples of the OUT distribution.

All experiments use $K = 32$ shadow models. Each shadow is trained for 3 epochs with the Adam optimiser, using the same hyperparameters as the target model (learning rate $2.5 \times 10^{-3}$, weight decay $1.5 \times 10^{-3}$, dropout 0.3). Three epochs keeps shadow training tractable at $K = 32$; whether additional epochs would change the attack success is not evaluated.

\subsection{Likelihood Ratio Score}
\label{ssec:mia_lr_score}

For each target subject $s$, the $K$ shadow models that had $s$ in their OUT set provide samples of the delta distribution under the hypothesis that $s$ was not in training. Let $\delta^{\text{out}}_k(s)$ denote the delta computed from shadow $k$ for subject $s$, where $s \notin$ shadow $k$'s training set. A Gaussian is fitted to these samples:
\begin{equation}
  \mu^{\text{out}}(s),\; \sigma^{\text{out}}(s)
  \;=\; \text{MLE fit of }
  \bigl\{\delta^{\text{out}}_k(s)\bigr\}_{k:\, s \notin \text{train}_k}.
\end{equation}
The offline LiRA score (Equation~\eqref{eq:lira_offline} from Chapter~\ref{ch:background}) is then:
\begin{equation}
  \text{score}(s)
  = \Phi\!\bigl(\delta_{\text{target}}(s);\,
    \mu^{\text{out}}(s),\, \sigma^{\text{out}}(s)\bigr),
\end{equation}
where $\delta_{\text{target}}(s)$ is the delta computed from the target model. A score close to 1 means the target model produces a higher delta for $s$ than the OUT shadow distribution predicts, indicating membership.

% ─────────────────────────────────────────────────────────────────────────────
\section{Threat Model Comparison}
\label{sec:mia_threat}

Three shadow initialisation strategies are compared. The grey-box variant assumes access to the target model checkpoint; the warm-start and cold-start variants both assume architecture-only knowledge and differ only in whether a shared surrogate is used to initialise the shadows. All three use the raw-delta formulation and the same $K = 32$ shadow budget.

\subsection{Grey-Box Attack}
\label{ssec:mia_grey}

The grey-box attacker has access to the target model checkpoint and initialises each of the $K = 32$ shadow LSTMs from those exact weights. Each shadow is then fine-tuned for 3 epochs on its 50\% member subset.

\subsection{Cold-Start Black-Box Attack}
\label{ssec:mia_cold}

The cold-start attacker knows only the model architecture. Each of the $K = 32$ shadow LSTMs is initialised independently from random Xavier weights and trained for 3 epochs on its 50\% member subset. There is no shared starting point between shadows.

\subsection{Warm-Start Black-Box Attack}
\label{ssec:mia_warm}

Like cold-start, the warm-start attacker knows only the model architecture. The difference is the initialisation strategy: rather than starting each shadow from independent random weights, the attacker first trains a single surrogate LSTM from random initialisation for 5 epochs on all available pairs, then uses those surrogate weights as the shared starting point for all $K$ shadow models. Each shadow is then fine-tuned for 3 epochs on its 50\% member subset. The motivation is to test whether this head start reduces effective training time per shadow without substantially reducing attack effectiveness.

\subsection{Results and Discussion}
\label{ssec:mia_threat_results}

Table~\ref{tab:threat_comparison} summarises the raw-delta AUC and TPR at 10\% FPR across all three threat models and four configurations.

\begin{table}[htbp]
\centering
\caption{LiRA results (raw-$\delta$) by threat model. AUC below 0.5 indicates
         inverted signal for grey-box and warm-start on whuGAIT (see Section~\ref{ssec:mia_threat_results}).}
\label{tab:threat_comparison}
\begin{tabular}{llrrrr}
\toprule
Config & Metric & Simple-$\delta$ & Grey & Warm & Cold \\
\midrule
\multirow{2}{*}{whuGAIT}
  & AUC          & \whuGAITSimplDAUC   & \whuGAITGreyLiraRawAUC  & \whuGAITWarmLiraRawAUC  & \whuGAITColdLiraRawAUC  \\
  & TPR at 10\% FPR  & \whuGAITSimplDTPRten & \whuGAITGreyLiraRawTPRten & \whuGAITWarmLiraRawTPRten & \whuGAITColdLiraRawTPRten \\
\multirow{2}{*}{UCI-HAR}
  & AUC          & \uciharSimplDAUC    & \uciharGreyLiraRawAUC   & \uciharWarmLiraRawAUC   & \uciharColdLiraRawAUC   \\
  & TPR at 10\% FPR  & \uciharSimplDTPRten  & \uciharGreyLiraRawTPRten  & \uciharWarmLiraRawTPRten  & \uciharColdLiraRawTPRten  \\
\multirow{2}{*}{WISDM}
  & AUC          & \wisdmSimplDAUC     & \wisdmGreyLiraRawAUC    & \wisdmWarmLiraRawAUC    & \wisdmColdLiraRawAUC    \\
  & TPR at 10\% FPR  & \wisdmSimplDTPRten   & \wisdmGreyLiraRawTPRten   & \wisdmWarmLiraRawTPRten   & \wisdmColdLiraRawTPRten   \\
\multirow{2}{*}{Combined}
  & AUC          & \combinedSimplDAUC  & \combinedGreyLiraRawAUC & \combinedWarmLiraRawAUC & \combinedColdLiraRawAUC \\
  & TPR at 10\% FPR  & \combinedSimplDTPRten & \combinedGreyLiraRawTPRten & \combinedWarmLiraRawTPRten & \combinedColdLiraRawTPRten \\
\bottomrule
\end{tabular}
\end{table}

Cold-start is the most consistent attack across datasets, achieving above-chance AUC in all three within-dataset configurations. The key reason is shadow diversity: when each shadow begins from an independent random initialisation, the $K$ OUT delta values for a given subject span a wider range of model behaviours, producing a more informative Gaussian fit for the LiRA score. Warm-start shadows all converge from the same surrogate, producing OUT distributions that are narrower and more correlated than cold-start's, failing to capture the full variability of the OUT hypothesis, sacrificing the diversity that makes cold-start effective.

Both grey-box and warm-start produce inverted signals on whuGAIT (AUC $= \whuGAITGreyLiraRawAUC$ and $\whuGAITWarmLiraRawAUC$ respectively, both below chance), while performing competitively on UCI-HAR and WISDM. Cold-start is the only variant that avoids inversion on whuGAIT. The common factor between grey-box and warm-start is that both initialise shadows from a pre-trained checkpoint (the target model or a surrogate) that has already seen all enrolled subjects. A plausible explanation is that 3 fine-tuning epochs are not sufficient to override this prior before the OUT subject is scored, inflating the OUT distribution and inverting the LiRA score. The cause is not confirmed; increasing the number of fine-tuning epochs per shadow is a natural direction to explore.

The combined configuration produces near-chance scores across all variants. With 159 training subjects spread across three datasets, the number of OUT samples per subject that satisfy the minimum validity threshold ($\geq 5$ shadow OUT deltas) is insufficient for reliable Gaussian fitting.

Figure~\ref{fig:threat_models} summarises the raw-delta AUC by threat model and dataset; Figure~\ref{fig:mia_roc} shows the corresponding ROC curves. The whuGAIT panel in Figure~\ref{fig:mia_roc} illustrates the signal inversion for grey-box and warm-start: both curves fall below the diagonal, indicating that the attack's scores are anti-correlated with membership. Cold-start is the only variant that remains above the diagonal on whuGAIT.

\begin{figure}[htbp]
  \centering
  \includegraphics[width=0.85\textwidth]{analysis/mia_nb05c_threat_models.png}
  \caption{LiRA AUC (raw-$\delta$) by threat model and dataset. The dashed line
           marks the random-classifier baseline (AUC $= 0.5$). Grey-box and
           warm-start fall below chance on whuGAIT; cold-start is above chance
           on all three within-dataset configurations.}
  \label{fig:threat_models}
\end{figure}

\begin{figure}[htbp]
  \centering
  \includegraphics[width=\textwidth]{analysis/mia_roc_comparison.png}
  \caption{MIA ROC curves for simple-$\delta$ and LiRA variants (grey-box
           raw, cold-start raw, cold-start logit-first) across all four
           dataset configurations. AUC values are shown in the legend. The
           whuGAIT panel shows grey-box falling below the diagonal (inverted
           signal), while cold-start sits well above it.}
  \label{fig:mia_roc}
\end{figure}

% ─────────────────────────────────────────────────────────────────────────────
\section{Cold-Start Delta Variants}
\label{sec:mia_delta_variants}

Having established cold-start as the most effective and robust LiRA threat model, three formulations of the delta score are compared under it. The raw delta in Equation~\eqref{eq:bg_delta} computes the difference of mean model outputs. Two alternatives transform the individual pair scores before aggregation, motivated by the same reasoning as the logit transform in the original LiRA for classification models.

\textbf{Raw} ($\delta_\text{raw}$): direct difference of mean scores,
$\bar{P}_\text{imp}(s) - \bar{P}_\text{gen}(s)$. Simple and interpretable, but the scale of the difference compresses near 0 and 1 where probabilities saturate.

\textbf{Logit-first} ($\delta_\text{lf}$): the logit transform is applied to
each individual pair score before computing the mean:
\begin{equation}
  \delta_\text{lf}(s)
  = \overline{\mathrm{logit}\!\left(P_\text{imp}(s)\right)}
  - \overline{\mathrm{logit}\!\left(P_\text{gen}(s)\right)}.
\end{equation}
This is the closest analogue to the original LiRA formulation: the logit is applied per sample before aggregation, stabilising scores at the extremes and giving greater weight to pairs where the model is highly confident.

\textbf{Mean-first} ($\delta_\text{mf}$): the logit is applied after computing
the mean:
\begin{equation}
  \delta_\text{mf}(s)
  = \mathrm{logit}\!\left(\bar{P}_\text{imp}(s)\right)
  - \mathrm{logit}\!\left(\bar{P}_\text{gen}(s)\right).
\end{equation}
This is more stable when the number of pairs per subject is small, because a single extreme pair score does not dominate the logit.

Table~\ref{tab:cold_variants} compares the three formulations under the cold-start threat model.

\begin{table}[htbp]
\centering
\caption{Cold-start LiRA results by delta variant. Simple-$\delta$ shown for reference.}
\label{tab:cold_variants}
\begin{tabular}{llrrrr}
\toprule
Config & Metric & Simple-$\delta$ & Cold-raw & Cold-lf & Cold-mf \\
\midrule
\multirow{2}{*}{whuGAIT}
  & AUC         & \whuGAITSimplDAUC   & \whuGAITColdLiraRawAUC  & \whuGAITColdLiraLfAUC  & \whuGAITColdLiraMfAUC  \\
  & TPR at 10\% FPR & \whuGAITSimplDTPRten & \whuGAITColdLiraRawTPRten & \whuGAITColdLiraLfTPRten & \whuGAITColdLiraMfTPRten \\
\multirow{2}{*}{UCI-HAR}
  & AUC         & \uciharSimplDAUC    & \uciharColdLiraRawAUC   & \uciharColdLiraLfAUC   & \uciharColdLiraMfAUC   \\
  & TPR at 10\% FPR & \uciharSimplDTPRten  & \uciharColdLiraRawTPRten  & \uciharColdLiraLfTPRten  & \uciharColdLiraMfTPRten  \\
\multirow{2}{*}{WISDM}
  & AUC         & \wisdmSimplDAUC     & \wisdmColdLiraRawAUC    & \wisdmColdLiraLfAUC    & \wisdmColdLiraMfAUC    \\
  & TPR at 10\% FPR & \wisdmSimplDTPRten   & \wisdmColdLiraRawTPRten   & \wisdmColdLiraLfTPRten   & \wisdmColdLiraMfTPRten   \\
\multirow{2}{*}{Combined}
  & AUC         & \combinedSimplDAUC  & \combinedColdLiraRawAUC & \combinedColdLiraLfAUC & \combinedColdLiraMfAUC \\
  & TPR at 10\% FPR & \combinedSimplDTPRten & \combinedColdLiraRawTPRten & \combinedColdLiraLfTPRten & \combinedColdLiraMfTPRten \\
\bottomrule
\end{tabular}
\end{table}

Raw and mean-first perform similarly across datasets. Logit-first is generally weaker, possibly because gait pair scores cluster near 0 and 1 more than classification confidences do, making the per-pair logit numerically unstable for a larger fraction of pairs. Figure~\ref{fig:cold_variants} shows the three formulations side by side.

\begin{figure}[htbp]
  \centering
  \includegraphics[width=0.85\textwidth]{analysis/mia_nb05c_cold_variants.png}
  \caption{Cold-start LiRA AUC by delta variant (raw, logit-first, mean-first)
           across all four configurations. Raw and mean-first track closely;
           logit-first underperforms on the combined run.}
  \label{fig:cold_variants}
\end{figure}

% ─────────────────────────────────────────────────────────────────────────────
\section{Signal Disentanglement}
\label{sec:mia_signal}

\subsection{Motivation and Design}
\label{ssec:mia_signal_design}

The experiments in the preceding sections train both the CNN encoder and the authentication LSTM on the same subject pool, so any measured membership signal could originate from either component. The CNN may have memorised the gait representations of its training subjects during classification pre-training; the LSTM may have fitted the pair-score distribution of its own training subjects during authentication training. These two contributions are inseparable in the standard pipeline.

To disentangle them, the pipeline is re-run on the \texttt{whuGAIT\_signal} variant, which partitions the 60 available subjects into three non-overlapping groups before any training:

\begin{itemize}
  \item \textbf{Group~A} (subjects 21–40): used to train the CNN encoder, but
    excluded from authentication training. Their gait representations are
    shaped by the CNN's classification objective.
  \item \textbf{Group~B} (subjects 41–60): excluded from CNN training, but
    used to train the authentication LSTM. The CNN has never seen their gait,
    so their CNN embeddings are generic; the LSTM has fitted their pair scores
    directly.
  \item \textbf{Group~C} (subjects 1–20): held out of both training phases.
    They serve as the non-member baseline.
\end{itemize}

Running the MIA attack with Group~B as the member class and Group~C as the non-member baseline isolates the authentication LSTM's memorisation signal, because the CNN treats both groups equally. Running it with Group~A as the member class and Group~C as the baseline isolates the CNN encoder's contribution, because the LSTM has not been trained on either group.

\subsection{Results}
\label{ssec:mia_signal_results}

Table~\ref{tab:signal_mia} reports the MIA AUC for Group~A (CNN signal) and Group~B (LSTM signal), each measured against the Group~C baseline, under the simple-delta and cold-start LiRA attacks.

\begin{table}[htbp]
\centering
\caption{Signal disentanglement MIA AUC. Group~B (LSTM members) and Group~A (CNN members).}
\label{tab:signal_mia}
\begin{tabular}{lrr}
\toprule
Attack & Group~B (LSTM) & Group~A (CNN) \\
\midrule
Simple-$\delta$    & \signalMiaAUCBvsC    & \signalMiaAUCAvsC \\
LiRA cold-start    & \signalMiaAUCBvsCCold & \signalMiaAUCAvsCCold \\
\bottomrule
\end{tabular}
\end{table}

Both groups produce above-chance AUC, confirming that each component independently leaks membership information. The LSTM signal (Group~B) is modestly stronger than the CNN signal (Group~A) under simple-delta (\signalMiaAUCBvsC{} vs \signalMiaAUCAvsC{}), but the gap is small. Under cold-start LiRA the two are essentially equal (\signalMiaAUCBvsCCold{} vs
\signalMiaAUCAvsCCold{}), suggesting that neither component dominates when the
attack is calibrated against an OUT distribution.

The authentication model's performance per group mirrors this pattern: the model achieves AUC of \signalAuthModelAUCB{} on Group~B pairs, \signalAuthModelAUCA{} on Group~A pairs, and \signalAuthModelAUCC{} on Group~C pairs. The fact that Group~A sits above the held-out baseline despite the LSTM never having trained on their pairs shows that the CNN encoder's prior representation carries a measurable advantage for the authentication head.

Figures~\ref{fig:signal_auc_bar} and~\ref{fig:signal_roc} show the results visually. The bar chart in Figure~\ref{fig:signal_auc_bar} shows mean percentile rank per group and attack variant: Groups~A and~B both sit above the random baseline (0.33), while Group~C falls below or near it. Figure~\ref{fig:signal_roc} shows the ROC curves for Group~B (LSTM signal) and Group~A (CNN signal) across all four attack variants, each using Group~C as the non-member baseline.

\begin{figure}[htbp]
  \centering
  \includegraphics[width=0.85\textwidth]{analysis/signal_mia_auc_per_group.png}
  \caption{Mean percentile rank of MIA attack score by subject group and attack
           variant. Higher rank means the attack places the group closer to the
           member end of the score distribution. Group~C (held-out) sits near
           the random baseline of 0.33.}
  \label{fig:signal_auc_bar}
\end{figure}

\begin{figure}[htbp]
  \centering
  \includegraphics[width=\textwidth]{analysis/signal_roc_per_group.png}
  \caption{MIA ROC curves for the signal disentanglement experiment. Left:
           Group~B (LSTM members) scored against the Group~C held-out baseline,
           isolating the authentication LSTM's contribution. Right: Group~A
           (CNN members) against the same baseline, isolating the CNN encoder's
           contribution.}
  \label{fig:signal_roc}
\end{figure}

\input{chapters/06_signal}
\input{chapters/07_evasion}
\input{chapters/08_conclusions}

% ── bibliography ──────────────────────────────────────────────────────────────
\bibliographystyle{plain}
\bibliography{references}

\end{document}
